\section{Critical estimates at a finite shape boundary}
\label{sec:critical}

We now consider the remaining nonisotropic case. Thus \(0<H<\infty\),
the functions \(f\) and \(q=f''/f\) are holomorphic in \(D_H\), and
\[
 M(s)=\max\{|f(s)|,|q(s)|^{1/2}\},\qquad
 \limsup_{k\to\infty}\frac{\log\max_x M(x+\ii(H-e^{-k}))}{k}=1.
\]
The bootstrap and no-trivial-profile results below use only this
critical growth condition and the results already proved. The
reductions in the last subsection have one explicit additional
hypothesis: the exclusion of the sequences in
\Cref{def:prohibited-sequence}. That exclusion is proved in
\Cref{thm:sequence-exclusion}, without using the angle budget or the
polynomial bounds that depend on that exclusion.

Choose \(k_0\ge 1\) large enough that \(d_{k_0}=e^{-k_0}<H/10\),
and write
\begin{equation}\label{critical:eq:means}
 \begin{split}
 d_k&=e^{-k},\\
 \mu_f(k)&=d_k\bigl(\E |f(x+\ii(H-d_k))|^2\bigr)^{1/2},\\
 \mu(k)&=d_k\bigl(\E(|f(x+\ii(H-d_k))|^2
                         +|q(x+\ii(H-d_k))|)\bigr)^{1/2}.
 \end{split}
\end{equation}
Unless a different probability is displayed, \(\E\) means integration
against \(\dd x/(2\pi)\) on the horizontal circle. In particular,
\[
 \|d_kM\|_2\le\mu(k)\le\sqrt{2}\|d_kM\|_2.
\]
The critical growth assumption gives, for every \(\eps>0\),
\begin{equation}\label{critical:eq:subexponential}
 \mu(k)+d_k\max_x M(x+\ii(H-d_k))
 \le C_\eps e^{\eps k}.
\end{equation}
Moreover, \Cref{lem:log-deficit}, applied on disks whose radii are fixed
small multiples of \(d_k\), gives
\begin{equation}\label{critical:eq:q-first-moment}
 d_k\E |q(x+\ii(H-d_k))|^{1/2}\le C(1+k).
\end{equation}
Indeed, on the slightly enlarged range of heights needed for those
disks, the critical assumption gives the logarithmic upper bound
\(\log|f|\le C(1+k)\). The horizontal mean of \(\log|f|\) is convex
in height and even under reflection in the real axis. It is therefore
bounded below by its finite real-axis mean. Averaging the disk deficits
in their horizontal centers gives \eqref{critical:eq:q-first-moment}.

Use the mean-normalized lift of \Cref{thm:shape-lift}, and put
\begin{equation}\label{critical:eq:P-and-deficit}
 P=-\Im V,\qquad W=f^2+8q,\qquad \Delta(t)=H-h_X(t).
\end{equation}
Here \(\mathfrak s(y,t)=y+t^2V(y)+O(t^4)\), and
\Cref{eq:V-coefficient} says that
\begin{equation}\label{critical:eq:gradient}
 V'=(W-\langle W\rangle)/120.
\end{equation}
Thus \(P\) is a real harmonic, periodic function, its horizontal mean
is zero on every row, and \(P=0\) on the real row. The number
\(\langle W\rangle\) is fixed throughout this Section. It follows from
\Cref{prop:width-superharmonic} and the normalized mean of the lift that
\(\Delta\) is a bounded nonnegative subharmonic function on a fixed
small parameter disk, with \(\Delta(0)=0\).

For real \(t\) in this disk and \(0<\Im y<h_X(t)\),
\begin{equation}\label{critical:eq:lower-lift}
 \Im\bigl(y-\mathfrak s(y,t)\bigr)\ge-\Delta(t).
\end{equation}
To verify this, the left side is harmonic and periodic in \(y\), and
is zero on the real row. On a row \(\Im y=h_X(t)-\eta\), the shape
constraint gives the lower bound \(-\Delta(t)-\eta\). Harmonic
comparison on that finite periodic strip, followed by \(\eta\downarrow0\),
proves \eqref{critical:eq:lower-lift}.

\subsection{Envelopes, tiles, and the one-sided estimate}

\begin{definition}\label{critical:def:envelope}
An envelope is a function \(N:[k_0,\infty)\to[1,\infty)\) for which
there are fixed constants \(C_N\ge1\), \(0\le\beta<1/2\), and
\(a>0\) such that
\begin{equation}\label{critical:eq:envelope}
 \frac{N(j)}{N(k)}\le C_N
 \begin{cases}
 e^{\beta(k-j)},& k_0\le j\le k,\\
 e^{a(j-k)},& j\ge k.
 \end{cases}
\end{equation}
\end{definition}

For the estimates below we impose the first-moment bound
\begin{equation}\label{critical:eq:first-moment}
 \E(d_kM)\le C_1N(k),\qquad k\ge k_0.
\end{equation}
When using a bootstrap, we also assume that
\(\sup_{k\ge k_0}\mu(k)/N(k)<\infty\); no bound on this finite
number is assumed. All conclusions below are uniform in it. In
particular, an envelope that eventually grows at any fixed positive
exponential rate has this finiteness property by
\eqref{critical:eq:subexponential}.

The constants in the ensuing estimates may depend on the table,
\(k_0,C_N,C_1,\beta,a\), and on specified fixed tile enlargements.
They do not depend on \(k\), on the finite bootstrap ratio, or on the
particular envelope in a family with these constants fixed.

For \(A\ge1\), a tile about \((k,x)\) will mean a set contained in
\begin{equation}\label{critical:eq:tile}
 \mathcal T_A(k,x)=
 \{s:\ e^{-A}d_k\le H-\Im s\le e^Ad_k,
          \ |\Re s-x|\le A(H-\Im s)\}.
\end{equation}
Horizontal distance is interpreted on a real lift and then
periodically. We intersect tiles with the upper half of the shape
strip. If a tile reaches a fixed lower region, we use the analytic
bounds there. In particular, all such fixed lower contributions are
harmless because \(N\ge1\).

We record the elementary averaging fact used repeatedly. For every
\(p>0\), the function \(M^p\) is subharmonic: it is the maximum of
\(|f|^p\) and \(|q|^{p/2}\). Covering a fixed tile by a bounded number
of disks of radius comparable to \(d_k\), and using the submean
inequality on slightly larger disks, bounds its supremum by
\(C d_k^{-2}\) times the integral on a fixed enlargement. Consequently,
\begin{equation}\label{critical:eq:tile-averaging}
 \E\sup_{s\in\mathcal T_A(k,x)}(d_kM(s))^p
 \le C_{A,p}\sup_{|j-k|\le C_A}\E(d_jM)^p,
\end{equation}
with the fixed lower-region convention just specified. One may replace
the supremum on the right by a bounded average of neighboring rows.
To obtain the displayed form, integrate the disk bound in \(x\);
each point of the enlargement is counted for a horizontal length
\(O(d_k)\), and its vertical range has length \(O(d_k)\).

Choose
\(\beta<\beta_1<1/2\), a sufficiently large fixed aperture \(L\),
a sufficiently small fixed \(c_0>0\), and a fixed
\(\delta_b\in(0,H/2)\). We take \(\delta_b\) sufficiently small that
the fixed enlargements of cone tiles used below, in particular in
\Cref{sign:large-moment}, have upper-boundary depth at most
\(e^{-k_0}\). This is possible after the fixed cone and tile apertures
have been chosen. Define
\begin{equation}\label{critical:eq:cone-cost}
 b(x,k)=1+
 \sup_{\substack{c_0d_k\le\delta\le\delta_b\\
                    |\Re s-x|\le L\delta\\
                    \delta=H-\Im s}}
 \frac{\delta M(s)}{N(k)(\delta/d_k)^{\beta_1}}.
\end{equation}
The constants \(L,c_0\) are chosen to cover the enlarged cones required
by \Cref{thm:cone}, with exponent \(1-\beta_1>1/2\), for a row point
and its needed fixed neighborhoods.

Splitting the cone into unit ranges of logarithmic depth proves
\begin{equation}\label{critical:eq:cone-moments}
 \E b(\cdot,k)\le C,
 \qquad
 \|b(\cdot,k)\|_p\le C_p(1+C_p^*)\quad(p\ge1),
\end{equation}
whenever \(\|d_jM\|_p\le C_p^*N(j)\) on all rows. Indeed, at a
coarser level \(j=k-h\), the normalized tile contribution is at most
\(C e^{-\beta_1h}N(j)/N(k)\le C e^{-(\beta_1-\beta)h}\)
times its normalized row norm. These factors are summable. There are
only a fixed number of finer levels because of \(c_0\). The
first inequality uses \eqref{critical:eq:first-moment}; the second
uses Minkowski's inequality and \eqref{critical:eq:tile-averaging}.

\begin{lemma}\label{critical:lem:one-sided}
Under \eqref{critical:eq:envelope} and
\eqref{critical:eq:first-moment}, there are constants \(c,z_0,C>0\)
such that, for real \(t\) and
\(z=|t|N(k)/d_k\le z_0\),
\begin{equation}\label{critical:eq:real-deficit}
 \Delta(t)\le C d_kz^\gamma,
 \qquad \gamma=\frac1{1+a}.
\end{equation}
Moreover, with
\begin{equation}\label{critical:eq:m-exponent}
 m=\frac{4(2-\gamma)}{4-\gamma}<2,
\end{equation}
one has
\begin{equation}\label{critical:eq:negative-P}
 \frac{d_kP^-(x+\ii(H-d_k))}{N(k)^2}
 \le C b(x,k)^m.
\end{equation}
Here \(P^-\) is the nonnegative negative part of \(P\).
\end{lemma}

\begin{proof}
At a fixed center the cone theorem gives a joint holomorphic vertical
passage at the row point for
\[
 |t|<r_{x,k}:=c\frac{d_k}{N(k)b(x,k)}.
\]
By decreasing \(c\), and using a fixed slightly shallower target in
the application of that theorem, it gives
\(|\mathfrak s-y|\le C d_k\) on this parameter disk and on a fixed
neighborhood of the row point. The passage is attached to the
actual lift from below and stays inside the shape strip. Evenness
in \(t\) and Cauchy's estimate therefore give
\begin{equation}\label{critical:eq:Taylor-error}
 \left|\mathfrak s(y,t)-y-t^2V(y)\right|
 \le C d_k(zb(x,k))^4
 \quad\text{when }|t|\le r_{x,k}/2.
\end{equation}
The fixed factors in \(c\) are included in the constant here.

By \eqref{critical:eq:cone-moments} there is a center with
\(b(x,k)\le C\). The real no-patch transfer from \Cref{thm:no-patch}
then implies
\begin{equation}\label{critical:eq:basic-width}
 \Delta(t)\le d_k
 \quad\text{for real }|t|\le c\,d_k/N(k).
\end{equation}
This holds for all sufficiently large \(k\); decreasing \(c\) handles
the remaining fixed compact range by joint analyticity near \(t=0\).

For \(j=k+h\), \eqref{critical:eq:envelope} gives
\[
 \frac{d_j}{N(j)}\ge C_N^{-1}
       \frac{d_k}{N(k)}e^{-(1+a)h}.
\]
For small \(z\), choose \(h\ge0\) so that
\(e^{-h}=C'z^{1/(1+a)}\), where the fixed constant \(C'\) is large
enough for \eqref{critical:eq:basic-width} to apply at level \(j\).
This proves \eqref{critical:eq:real-deficit}. Rounding \(h\) to an
integer, if desired, changes only a fixed constant.

For the negative part, take
\(z=c'b^{-4/(4-\gamma)}\), with \(c'>0\) fixed sufficiently small.
Since \(b\ge1\), this makes both \(zb\) and
\(\Delta(t)/d_k\le Cz^\gamma\) small. In particular the row lies
strictly inside the actual strip. Its actual lift agrees with the
vertical passage in \eqref{critical:eq:Taylor-error}. Combining that
estimate with \eqref{critical:eq:lower-lift} gives
\[
 \frac{d_kP^-}{N(k)^2}
 \le C\left(z^{\gamma-2}+z^2b^4\right).
\]
The two powers of \(b\) coincide and equal the exponent in
\eqref{critical:eq:m-exponent}. Since \(\gamma>0\), that exponent
is strictly less than two.
\end{proof}

\subsection{The quadratic bootstrap and its moment improvement}

\begin{lemma}\label{critical:lem:disk-bound}
On an inner tile, with a fixed enlargement on the right,
\begin{equation}\label{critical:eq:disk-bound}
 (d_kM)^2\le C\left(1+\sup d_k^2|W|\right).
\end{equation}
\end{lemma}

\begin{proof}
It is enough to prove the assertion on nested disks. In disk
coordinates put \(g=d_kf\) and \(\widetilde W=d_k^2W\). Then
\[
 g''=\frac18g(\widetilde W-g^2).
\]
Suppose that no universal bound
\(\sup_{\mathbb D}|g|\le C(1+\sup_{2\mathbb D}|\widetilde W|)^{1/2}\)
holds. For a counterexample sequence choose in an intermediate disk a
point maximizing \(|g(z)|\) times its distance to that disk's boundary.
Write \(m_n\) for the value of \(|g|\) there and \(r_n\) for that
distance. Failure of the proposed estimate implies
\(r_nm_n\to\infty\) and
\(\sup|\widetilde W|/m_n^2\to0\). On the disk of radius \(r_n/2\)
about the chosen point, maximality gives \(|g|\le2m_n\).
Rescale space by \(m_n^{-1}\) and divide the function by \(m_n\).
After passage to a subsequence, the resulting functions converge on
every compact subset of \(\C\) to an entire function \(G\) with
\(|G|\le2\), \(|G(0)|=1\), and
\(G''=-G^3/8\). Liouville's theorem makes \(G\) constant, which
contradicts this equation and \(|G(0)|=1\). This bounds \(f\).
The identity \(8q=W-f^2\) bounds \(q\) on the same inner disk.
A finite disk covering proves \eqref{critical:eq:disk-bound}.
\end{proof}

\begin{theorem}[Critical moment bootstrap]\label{thm:critical-bootstrap}
Let \(N\) be an envelope as in \Cref{critical:def:envelope}, assume
the first-moment bound \eqref{critical:eq:first-moment}, and suppose
\(\sup_{k\ge k_0}\mu(k)/N(k)<\infty\). Then
\begin{equation}\label{critical:eq:L2-bootstrap}
 \mu(k)\le C N(k),\qquad
 \E\frac{d_k|P|}{N(k)^2}\le C.
\end{equation}
There is \(\eta>0\), depending only on the fixed envelope constants,
such that
\begin{equation}\label{critical:eq:higher-moment}
 \left\|\frac{d_kM}{N(k)}\right\|_{2+\eta}\le C.
\end{equation}
The estimates for \(M\) hold also with suprema on every fixed tile.
If \(N\) is nondecreasing, then in addition
\begin{equation}\label{critical:eq:global-envelope}
 \sup_x d_kM\le C(1+k)N(k),\qquad
 \sup_x\frac{d_k|P|}{N(k)^2}\le C(1+k)^2.
\end{equation}
\end{theorem}

\begin{proof}[Proof of the quadratic and higher-moment assertions]
Let \(C_* =\sup_{j\ge k_0}\mu(j)/N(j)<\infty\).
By \eqref{critical:eq:cone-moments},
\(\|b(\cdot,k)\|_2\le C(1+C_*)\). Hence
\eqref{critical:eq:negative-P} gives
\[
 \E\frac{d_kP^-}{N(k)^2}\le C(1+C_*^m).
\]
The horizontal mean of \(P\) is zero, so its positive and negative
parts have equal integrals. Thus the same bound, with a factor two,
holds for \(|P|\).

Interior harmonic estimates bound the supremum of \(|\nabla P|\)
on a tile by \(C d_k^{-3}\) times the integral of \(|P|\) on a fixed
enlargement. Averaging in the horizontal center and using the bounds
on neighboring rows gives
\[
 \E\sup_{\mathcal T_A(k,x)}
       \frac{d_k^2|W|}{N(k)^2}
 \le C_A(1+C_*^m).
\]
Here \eqref{critical:eq:gradient} accounts for the fixed additive
constant \(\langle W\rangle\), which is absorbed because \(N\ge1\).
The disk bound now yields
\(C_*^2\le C(1+C_*^m)\). Since \(m<2\), this bounds \(C_*\)
by a fixed constant and proves \eqref{critical:eq:L2-bootstrap},
including the tile version of its first assertion.

We give the concentration estimate that improves the moment. Put
\(p=2/m>1\) and \(\vartheta=1-1/p>0\). The bound for \(b\) and
\eqref{critical:eq:negative-P} give, on all rows in any fixed range of
depth ratios,
\begin{equation}\label{critical:eq:negative-Lp}
 \left\|\frac{d_kP^-}{N(k)^2}\right\|_p\le C.
\end{equation}
Partition the horizontal circle into intervals of lengths comparable
to \(d_k\). Select a nonempty collection of fractional total length
\(\sigma\), and let \(E\) be its enlargement by a fixed multiple of
\(d_k\). In particular \(\sigma\ge c d_k\), and
\(|E|/(2\pi)\le C\sigma\). We estimate \(P\) on \(E\) on a row
of depth \(d'\) comparable to \(d_k\).

Choose \(R\ge R_0\), with \(R_0\) fixed sufficiently large, and
suppose \(Rd_k\) is small. Harmonic extension from that row to depth
\(Rd_k\), with zero data on the real row, is convolution with a
positive periodic strip kernel \(K\) of mass at most one. Before
periodization its density is
\begin{equation}\label{critical:eq:strip-kernel}
 \frac1{2h'}
 \frac{\sin(\pi l/h')}
      {\cosh(\pi u/h')-\cos(\pi l/h')},
 \qquad h'=H-d',\quad l=Rd_k-d'.
\end{equation}
It is at least \(c/(Rd_k)\) when \(|u|\le Rd_k\). Let \(E_R\)
be the enlargement of \(E\) by a sufficiently large multiple of
\(Rd_k\); it has fractional length at most
\(C\min\{1,R\sigma\}\). Integrating first in the destination
variable shows that a fixed fraction of the positive mass of \(P\)
on \(E\) contributes to \(\int_{E_R}K*P^+\). As
\(K*P^+=P_{\mathrm{coarse}}+K*P^-\), we obtain
\begin{equation}\label{critical:eq:selected-P-plus}
 \int_E P^+\,\frac{\dd x}{2\pi}
 \le C\left(\int |P_{\mathrm{coarse}}|\,\frac{\dd x}{2\pi}
                  +\int_{E_R}K*P^-\,\frac{\dd x}{2\pi}\right).
\end{equation}
The normalized first term is at most \(CR^{2\beta-1}\): at the
coarse level \(j=k-\log R\), use
\(N(j)\le C R^\beta N(k)\) and \(d_j=Rd_k\) in
\eqref{critical:eq:L2-bootstrap}. The convolution is an \(L^p\)
contraction, so Holder's inequality and
\eqref{critical:eq:negative-Lp} bound the second normalized term by
\(C(R\sigma)^\vartheta\). The negative mass directly on \(E\)
has the same bound. Consequently
\begin{equation}\label{critical:eq:concentration-P}
 \int_E\frac{d_k|P(x+\ii(H-d'))|}{N(k)^2}
                \,\frac{\dd x}{2\pi}
 \le C\left(R^{2\beta-1}+(R\sigma)^\vartheta\right).
\end{equation}

The interior harmonic estimates used above, followed by
\Cref{critical:lem:disk-bound}, now give
\begin{equation}\label{critical:eq:concentration-M}
 \sum_{I\ \mathrm{selected}} |I|
       \sup_{x\in I}\sup_{s\in\mathcal T_A(k,x)}
                  \frac{d_k^2M(s)^2}{N(k)^2}
 \le C_A\left(\sigma+R^{2\beta-1}+(R\sigma)^\vartheta\right).
\end{equation}
To see that no concentration is lost in this passage, enlarge the
selected intervals by the fixed horizontal margins needed for the
harmonic estimates and integrate those estimates over the fixed
range of neighboring depths. The enlarged intervals have bounded
overlap and still total length \(O(\sigma)\), so
\eqref{critical:eq:concentration-P} applies at every such depth.

Take \(R=R_0\sigma^{-1/2}\). This is allowed for all sufficiently
large \(k\), since \(\sigma\ge c d_k\) gives
\(Rd_k\le C\sqrt{d_k}\). The remaining compact range of \(k\)
is harmless. There is therefore \(0<\zeta<1\), depending only on
\(\beta\) and \(p\), such that the right side of
\eqref{critical:eq:concentration-M} is at most \(C_A\sigma^\zeta\).
For example one may take any positive \(\zeta\) no larger than
\(\min\{(1-2\beta)/2,\vartheta/2\}\).

Let \(A_k(I)\) be the supremum appearing in the sum in
\eqref{critical:eq:concentration-M}. Selecting the cells for which
\(A_k(I)>b_0\ge1\) gives
\[
 \Prob\{A_k>b_0\}\le C b_0^{-1/(1-\zeta)}.
\]
The exponent is greater than one. Integrating this distribution
bound shows that \(A_k\) has a uniformly bounded moment of some
order \(1+\eta/2>1\). This proves
\eqref{critical:eq:higher-moment}, including its fixed-tile versions.
\end{proof}
\begin{proof}[Proof of the global assertions in
  \Cref{thm:critical-bootstrap}]
Suppose \(N\) is nondecreasing. We first prove that, for a fixed
small \(c_2>0\),
\begin{equation}\label{critical:eq:complex-deficit}
 \Delta(t)\le d_k/4\qquad
 \left(|t|\le \rho_k:=\frac{c_2d_k}{(1+k)N(k)}\right).
\end{equation}
Choose a fixed \(\eps_0>0\) so small that
\eqref{critical:eq:real-deficit} bounds \(\Delta\) by \(d_k/16\)
on the real interval \(|u|\le u_0:=\eps_0d_k/N(k)\).
For a fixed small parameter half-disk, harmonic upper comparison for
the bounded subharmonic function \(\Delta\) gives, when
\(|t|\le u_0/2\), an upper bound consisting of this near-diameter
contribution, an \(O(|t|)\) contribution from the curved boundary,
and
\begin{equation}\label{critical:eq:diameter-integral}
 C|t|\int_{u_0}^{r_0}
         \max\{\Delta(u),\Delta(-u)\}\,\frac{\dd u}{u^2}.
\end{equation}
Here \(r_0>0\) is fixed. The kernel bound follows directly from the
upper-half-plane Poisson kernel on the diameter; the harmonic measure
of the distant curved arc is \(O(|t|)\). Upper semicontinuous
boundary data cause no difficulty: apply the comparison first to
continuous majorants and decrease them to the given data. The lower
half-disk is handled in the same way.

We claim that the integral in \eqref{critical:eq:diameter-integral}
is at most \(C_{\eps_0}(1+k)N(k)\). For earlier levels \(j\le k\),
monotonicity gives \(N(j)\le N(k)\), so
\eqref{critical:eq:basic-width} implies
\(\Delta(u)\le d_j\) up to \(|u|\le c d_j/N(k)\).
Divide the interval of integration into the corresponding ranges
for \(j=k,k-1,\ldots\), down to a fixed level. On each such range
\(\Delta(u)\le C d_j\) and its \(u^{-2}\) integral is at most
\(C N(k)\). The first range costs \(C_{\eps_0}N(k)\).
Beyond the last one the boundedness of \(\Delta\) costs at most
\(C N(k)\). There are \(O(1+k)\) ranges, proving the claim.
Choosing \(c_2\) small now proves
\eqref{critical:eq:complex-deficit}.

For \(y\) at depth comparable to \(d_k\), with a fixed neighborhood
of that size, the actual lift is consequently holomorphic for
\(|t|<\rho_k\) and satisfies
\(\Im(\mathfrak s(y,t)-y)\le C d_k\). It vanishes at \(t=0\).
The coefficient estimate for an analytic function with positive real
part, applied to \(C d_k+\ii(\mathfrak s-y)\), gives
\[
 |V(y)|\le C d_k\rho_k^{-2}.
\]
This elementary estimate also follows by integrating the nonnegative
real part on concentric circles and taking its second Fourier
coefficient. It proves the asserted bound for \(P\).
Cauchy's estimate in \(y\), followed by
\eqref{critical:eq:gradient}, gives \(|W|\le C\rho_k^{-2}+C\)
on a slightly smaller tile. The disk bound therefore gives
\(d_kM\le C d_k/\rho_k+C\), which is the other assertion in
\eqref{critical:eq:global-envelope}.
\end{proof}

\subsection{A weighted version with an ancestor row}

We will need to concentrate horizontal centers near a selected
ancestor, while retaining estimates at much finer depths. For
\(0<d_a\ll1\), let \(K_{d_a}\) be a normalized periodic probability
kernel whose density, relative to \(\dd x/(2\pi)\), is uniformly
comparable to
\begin{equation}\label{critical:eq:ancestor-kernel}
 \frac{d_a}{d_a^2+|x|_{\mathrm{per}}^2}.
\end{equation}
The normalization factor in this comparison is bounded above and
below independently of \(d_a\). Given any probability \(\omega\)
on the circle, put \(\nu=K_{d_a}*\omega\). The following elementary
properties will be used.

First, for every fixed \(A\), the density of \(\nu\) changes by at
most a fixed factor under translations of length at most \(A d_a\).
For the kernel this follows from the triangle inequality and
\[
 d_a^2+|u+h|_{\mathrm{per}}^2
 \le C_A(d_a^2+|u|_{\mathrm{per}}^2)
 \quad (|h|\le A d_a),
\]
and convolution preserves the comparison. Second, if
\(d_j<d_a/2\), the unnormalized upper-boundary harmonic-measure law
from the ancestor point \(x+\ii(H-d_a)\) to the row at depth \(d_j\),
with real bottom, is uniformly comparable to the translated kernel
\(K_{d_a}\). Indeed in \eqref{critical:eq:strip-kernel} the distance
from the upper row lies between \(d_a/2\) and \(d_a\). On a
fundamental horizontal interval the denominator is comparable to
\(d_a^2+|u|^2\), with constants depending only on \(H\);
periodization contributes a comparable bounded tail. Mixing over
\(\omega\) gives an unnormalized law comparable to \(\nu\).

\begin{proposition}\label{prop:weighted-bootstrap}
Let \(N\) satisfy the hypotheses of \Cref{thm:critical-bootstrap},
with fixed constants, and let \(d_a=e^{-k_a}\) be a sufficiently
small ancestor depth. Let \(\omega,\nu\) be as above, with fixed
kernel comparison constants. Suppose that for every fixed tile
enlargement needed below,
\begin{equation}\label{critical:eq:weighted-coarse}
 \left\|\sup_{s\in\mathcal T_A(j,x)}d_jM(s)\right\|_{L^2(\nu)}
 \le C_A N(j),\qquad k_0\le j\le k_a+3,
\end{equation}
and that
\begin{equation}\label{critical:eq:weighted-anchor}
 \int |P(x+\ii(H-d_a))|\,\dd\omega(x)
 \le C\frac{N(k_a)^2}{d_a}.
\end{equation}
Then on all finer levels, and with every fixed tile enlargement,
\begin{equation}\label{critical:eq:weighted-conclusion}
 \left\|\sup_{s\in\mathcal T_A(j,x)}d_jM(s)\right\|_{L^2(\nu)}
 \le C'_A N(j).
\end{equation}
The constants are independent of \(k_a\), of \(\omega\), and of the
finite bootstrap ratio. No higher weighted moment is asserted.
\end{proposition}

\begin{proof}
It suffices first to bound row norms. Let \(C_*\) be their supremum
ratio to \(N\) on the fine levels, enlarged to include the fixed
coarse constants. This number is finite. For example, the density of
\(\nu\) is bounded for each fixed ancestor, so finiteness follows
from the unweighted finite-ratio hypothesis. This preliminary bound
may depend on \(d_a\); none of the estimates below will depend on it.

Tile averaging on levels whose depths are at most a fixed multiple
of \(d_a\) is valid for \(\nu\), by its translation-comparison
property. In the logarithmic summation defining \(b\), use these
fine tile bounds on that part of the cone, and use
\eqref{critical:eq:weighted-coarse} on the coarser part. This gives
\begin{equation}\label{critical:eq:weighted-b}
 \|b(\cdot,j)\|_{L^2(\nu)}\le C(1+C_*).
\end{equation}
Fixed ranges immediately adjacent to the coarse cutoff are already
covered by the assumed fixed tile enlargements at levels at most
\(k_a+3\). In particular, choosing larger fixed interior-estimate
neighborhoods does not create a gap at the cutoff.

For \(d_j<d_a/2\), let \(\kappa_j\) be the mixed unnormalized
harmonic-measure law from \(\omega\) to the row at depth \(d_j\).
Since \(P=0\) on the real bottom,
\[
 \int P(x+\ii(H-d_j))\,\dd\kappa_j(x)
   =\int P(x+\ii(H-d_a))\,\dd\omega(x).
\]
The laws \(\kappa_j\) and \(\nu\) are comparable. Moreover the
absolute value of the right side, in fine-row normalization, is at
most
\begin{equation}\label{critical:eq:anchor-decay}
 C\frac{d_j}{d_a}\frac{N(k_a)^2}{N(j)^2}
 \le C e^{-(1-2\beta)(j-k_a)}\le C.
\end{equation}
The negative-part estimate \eqref{critical:eq:negative-P}, integrated
using \eqref{critical:eq:weighted-b}, bounds the normalized negative
integral by \(C(1+C_*^m)\). Taking the signed integral under
\(\kappa_j\) first, and then using comparison of the two positive
measures, therefore gives
\[
 \int\frac{d_j|P(x+\ii(H-d_j))|}{N(j)^2}\,\dd\nu(x)
 \le C(1+C_*^m).
\]
The strict inequality \(2\beta<1\) gives uniform decay in
\eqref{critical:eq:anchor-decay}.

Weighted interior harmonic estimates now follow by the same local
averaging argument, since the required horizontal translations are
of size \(O(d_j)\). Combining them with
\Cref{critical:lem:disk-bound} gives
\(C_*^2\le C(1+C_*^m)\), away from the already controlled fixed
cutoff range. This bounds \(C_*\) uniformly. One more application of
weighted tile averaging proves \eqref{critical:eq:weighted-conclusion}.
\end{proof}

\subsection{A vanishing profile at one center is impossible}

The next result is deliberately stated for individual centers. A
single quiet cone forces a global real-parameter width estimate, and
the latter can then be tested against unweighted horizontal moments.

\begin{theorem}[Exclusion of vanishing profiles at individual centers]\label{thm:no-trivial-profile}
Consider sequences of levels \(k_n\to\infty\), scales
\(d_n=e^{-k_n}\), and amplitudes \(B_n=\mu(k_n)\to\infty\).
Suppose there are envelopes \(N_n\) satisfying the hypotheses of
\Cref{thm:critical-bootstrap} with constants independent of \(n\),
and \(N_n(k_n)\le C B_n\).

For a fixed \(\alpha\in(1/2,1)\) and a sufficiently large fixed
aperture \(L\), there cannot be centers \(x_n\) such that, for every
fixed \(c>0\),
\begin{equation}\label{critical:eq:quiet-cone}
 \sup_{\substack{c d_n\le\delta\le\delta_b\\
                  |\Re s-x_n|\le L\delta\\
                  \delta=H-\Im s}}
 \frac{d_nM(s)}{B_n}\left(\frac{\delta}{d_n}\right)^\alpha
 \longrightarrow0.
\end{equation}
In particular, suppose the profiles
\begin{equation}\label{critical:eq:profiles}
 F_n(S)=\frac{d_n}{B_n}f(x_n+\ii H+d_nS),\qquad
 Q_n(S)=\left(\frac{d_n}{B_n}\right)^2q(x_n+\ii H+d_nS)
\end{equation}
converge to \((0,0)\) on compact subsets of the lower half-plane.
They cannot also satisfy, for some fixed \(\alpha_0>1/2\) and an
aperture \(L\) sufficiently large for a chosen exponent
\(1/2<\alpha<\min(1,\alpha_0)\),
\begin{equation}\label{critical:eq:bounded-coarse-cone}
 \frac{d_n}{B_n}M(s)\le C\left(\frac{d_n}{\delta}\right)^{\alpha_0}
 \quad(d_n\le\delta\le\delta_b,
             \ |\Re s-x_n|\le L\delta).
\end{equation}
Every compact holomorphic limit of the pair in
\eqref{critical:eq:profiles} is pure: its product \(FQ\) vanishes
identically.
\end{theorem}

\begin{proof}
Suppose first that \eqref{critical:eq:quiet-cone} holds. For any fixed
\(c,T_0>0\), apply \Cref{thm:cone} at a depth which is a fixed
small multiple of \(c d_n\), and with parameter radius
\(T_0d_n/B_n\). The ratio of this radius to the target depth tends
to zero. Thus its fixed \(K=1\) regime suffices, and its aperture
can be chosen once, independently of \(c,T_0\). The weighted
smallness in \eqref{critical:eq:quiet-cone} supplies its small
parameter. The real no-patch transfer gives
\[
 \Delta(Td_n/B_n)\le c d_n
 \quad (|T|\le T_0)
\]
for all sufficiently large \(n\). Since \(c>0\) is arbitrary,
\begin{equation}\label{critical:eq:vanishing-width}
 \sup_{|T|\le T_0}\frac{\Delta(Td_n/B_n)}{d_n}\longrightarrow0
 \qquad(T\text{ real}).
\end{equation}

We show that this width estimate alone contradicts
\(B_n=\mu(k_n)\). Let \(j-k_n\) vary in an arbitrary fixed compact
interval. The envelope inequalities imply \(N_n(j)\le C B_n\),
with a constant uniform on that interval. By the quadratic bootstrap
and \eqref{critical:eq:cone-moments}, the cone costs \(b(x,j)\)
have uniformly bounded second moments. Also
\eqref{critical:eq:negative-P} gives
\begin{equation}\label{critical:eq:uniform-integrability}
 \frac{d_nP^-(x+\ii(H-d_j))}{B_n^2}\le C b(x,j)^m,
 \qquad m<2.
\end{equation}
These variables are uniformly integrable: their \(2/m\)-moments
are bounded, and \(2/m>1\).

Fix a cutoff \(A>1\) and work where \(b(x,j)\le A\). Take
\(t=T d_n/B_n\) with a fixed nonzero \(T\) small enough for the
Taylor estimate at all these points. By
\eqref{critical:eq:vanishing-width}, the relevant rows lie in the
actual strip for large \(n\). Combining
\eqref{critical:eq:lower-lift} and
\eqref{critical:eq:Taylor-error}, now in this normalization, gives
\[
 \frac{d_nP^-}{B_n^2}
 \le C\frac{\Delta(Td_n/B_n)}{T^2d_n}+C T^2 A^4.
\]
Given a positive threshold, first choose \(A\) so that the complement
has small probability, then choose \(|T|\) so that the second term
is below that threshold, and finally let \(n\to\infty\). This
proves convergence to zero in probability. The uniform integrability
in \eqref{critical:eq:uniform-integrability} gives convergence in
mean, uniformly on every fixed range of \(j-k_n\).

The signed horizontal mean of \(P\) is zero, so the same conclusion
holds for \(d_n|P|/B_n^2\). Interior harmonic estimates, using a
fixed range of neighboring depths, imply
\[
 \E\sup_{\text{fixed tile}}
       \frac{d_n^2|W|}{B_n^2}\longrightarrow0.
\]
The disk bound and \(B_n\to\infty\) now imply
\[
 \frac{\mu(k_n)^2}{B_n^2}\longrightarrow0,
\]
contradicting its value one.

For the stated consequence, choose
\(1/2<\alpha<\min\{1,\alpha_0\}\). On the part of the cone with
\(\delta/d_n\ge R\), \eqref{critical:eq:bounded-coarse-cone}
bounds the expression in \eqref{critical:eq:quiet-cone} by
\(C R^{\alpha-\alpha_0}\). On each remaining compact part with
\(c\le\delta/d_n\le R\), the assumed profile convergence makes
it tend to zero. Letting \(R\to\infty\) proves
\eqref{critical:eq:quiet-cone}, which was just excluded. In fact
compact convergence within this fixed cone is enough for the argument.

Finally the scaled differential identity is
\[
 \partial_S^2F_n=B_n^2Q_nF_n.
\]
Compact holomorphic convergence bounds its left side. Dividing by
\(B_n^2\) and passing to the limit gives \(FQ=0\). On the
connected lower half-plane, a nonzero analytic factor then forces
the other factor to vanish identically.
\end{proof}

\subsection{The sequence to be excluded and the angle-budget reduction}

We specify the input to the next Section before deriving the
consequences of its exclusion. For a probability \(\nu_n\) on
horizontal centers, convergence in probability uniformly on a compact
set means that the probability of the supremum exceeding each fixed
positive threshold tends to zero.

\begin{definition}\label{def:prohibited-sequence}
A prohibited sequence consists of the following data in the critical
setting of this Section:
\begin{enumerate}[label=\textup{(\roman*)}]
 \item Levels \(k_n\to\infty\), amplitudes
 \(B_n=\mu(k_n)\to\infty\), and numbers \(\eps_n\downarrow0\)
 such that
 \begin{equation}\label{critical:eq:primary-normality}
  \mu(j)\le B_n e^{\eps_n|j-k_n|}
       \qquad(j\ge k_0).
 \end{equation}
 \item Ancestors \(k_{a,n}=k_n-s_n\), with both
 \(k_{a,n}\to\infty\) and \(s_n\to\infty\), and probabilities
 \(\omega_n\) and
 \(\nu_n=K_{d_{a,n}}*\omega_n\), where
 \(d_{a,n}=e^{-k_{a,n}}\) and the kernels obey
 \eqref{critical:eq:ancestor-kernel} with comparison constants
 independent of \(n\).
 \item For every fixed \(0<\beta<1/2\) and every fixed tile
 enlargement \(A\), there is a constant \(C_{\beta,A}\) such that,
 eventually in \(n\), simultaneously for all
 \(k_0\le j\le k_{a,n}+3\),
 \begin{equation}\label{critical:eq:prohibited-coarse}
  \left\|\sup_{s\in\mathcal T_A(j,x)}d_jM(s)\right\|_{L^2(\nu_n)}
   \le C_{\beta,A} B_n e^{\beta(k_n-j)}.
 \end{equation}
 For every such \(\beta\) there is also a constant \(C_\beta\)
 such that, eventually,
 \begin{equation}\label{critical:eq:prohibited-anchor}
  \int |P(x+\ii(H-d_{a,n}))|\,\dd\omega_n(x)
   \le C_\beta
       \frac{B_n^2e^{2\beta(k_n-k_{a,n})}}{d_{a,n}}.
 \end{equation}
 The constants and the threshold for ``eventually'' may depend on
 the indicated fixed choices, but not on \(j\) or \(n\).
 \item With \(d_n=e^{-k_n}\), for every compact subset
 \(\mathcal K\) of \(\{\Im S<0\}\),
 \begin{equation}\label{critical:eq:prohibited-Q}
  \sup_{S\in\mathcal K}
  \left|\left(\frac{d_n}{B_n}\right)^2
          q(x+\ii H+d_nS)\right|
  \longrightarrow0
  \quad\text{in }\nu_n\text{-probability}.
 \end{equation}
\end{enumerate}
\end{definition}

For fixed \(\beta,a>0\), with \(\beta<1/2\), primary normality
controls the unweighted moments by an envelope of the form
\[
 N_n(j)=B_n
   \exp\bigl(\beta(k_n-j)_++a(j-k_n)_+\bigr).
\]
Thus \Cref{prop:weighted-bootstrap} applies to this envelope using
\eqref{critical:eq:prohibited-coarse} and
\eqref{critical:eq:prohibited-anchor}. This observation uses none of
the conditional conclusions below. It is the reason the ancestors
are included in the definition.

We will also use an elementary way to propagate a logarithmic
deficit. If \(U\ge0\) is superharmonic and two small disks have
comparable radii, nearby centers, and a sufficiently enlarged common
disk contained in its domain, then
\begin{equation}\label{critical:eq:deficit-step}
 \operatorname{av}_{D_2}U\le C\operatorname{av}_{D_1}U.
\end{equation}
Indeed, for each point of \(D_1\), a disk centered there and of a
fixed multiple of the common radius contains \(D_2\) and remains
inside the domain. Positivity and the supermean inequality bound the
average on \(D_2\) by a constant times the value at that point;
integration over \(D_1\) proves the claim. Starting with a disk
centered at a point where \(U\) is finite, the first average is at
most that value. A chain of \(m\) such steps therefore costs at most
\(C^m\). In the upper strip, one can connect points with logarithmic
depth separation \(s\), and with a horizontal displacement bounded
by a fixed multiple of the larger depth, by \(O(1+s)\) disks of
radii small fixed multiples of their depths. All enlarged disks can
be kept in a fixed slightly wider range of logarithmic depths.
This argument applies to \(K-\log|f|\), including its logarithmic
poles at zeros; it uses disk averages, not a pointwise Harnack
inequality for superharmonic functions.

\begin{proposition}\label{critical:prop:power-saving}
Assume that no sequence in \Cref{def:prohibited-sequence} exists.
Then there are \(C>0\) and \(\eps_1\in(0,1)\) such that
\begin{equation}\label{critical:eq:mu-f-saving}
 \mu_f(k)\le C(1+k)^{1-\eps_1}\qquad(k\ge k_0).
\end{equation}
\end{proposition}

\begin{proof}
Suppose otherwise. We first choose levels \(\kappa_n\to\infty\)
and values \(b_n=\mu_f(\kappa_n)\) such that
\begin{equation}\label{critical:eq:f-record}
 \begin{split}
 b_n&\ge\kappa_n^{1-o(1)},\\
 \mu_f(j)&\le b_n\quad(k_0\le j\le\kappa_n),\\
 \mu_f(j)&\le b_n e^{\eta_n(j-\kappa_n)}
       \quad(j\ge\kappa_n),\qquad \eta_n\to0.
 \end{split}
\end{equation}
Here is an explicit selection. Failure of
\eqref{critical:eq:mu-f-saving} provides arbitrarily large levels
with \(\mu_f\ge k^{1-1/n}\). Replace such a level by a maximum
of \(\mu_f\) over all earlier levels; denote the resulting record
level by \(u_n\). The record values tend to infinity, so
\(u_n\to\infty\), and they still exceed \(u_n^{1-1/n}\).
Choose \(\eta_n\to0\) with \(\eta_n\ge1/u_n\). Maximize
\(\mu_f(j)e^{-\eta_n(j-u_n)}\) on \(j\ge u_n\), and call a
maximizer \(\kappa_n\). It exists by
\eqref{critical:eq:subexponential}. It is still a record among all
earlier levels. Its value satisfies
\[
 b_n\ge u_n^{1-1/n}e^{\eta_n(\kappa_n-u_n)}
       \ge\kappa_n^{1-1/n},
\]
where we used \(\log(\kappa_n/u_n)\le(\kappa_n-u_n)/u_n\).
The maximizing property gives the last inequality in
\eqref{critical:eq:f-record}.

For each fixed \(\eta>0\), and all sufficiently large \(n\), set
\[
 N_{n,\eta}(j)=(b_n+\kappa_n)
                   e^{\eta(j-\kappa_n)_+}.
\]
This nondecreasing envelope controls the first moment uniformly in
\(n\). For \(|f|\) use \eqref{critical:eq:f-record}; for
\(|q|^{1/2}\) use \eqref{critical:eq:q-first-moment} and
\(1+j\le C_\eta(1+\kappa_n)e^{\eta(j-\kappa_n)}\) on the
future side. Its finite-ratio hypothesis follows from its exponential
tail. \Cref{thm:critical-bootstrap} consequently gives
\begin{equation}\label{critical:eq:record-envelope}
 \mu(j)\le C_\eta(b_n+\kappa_n)e^{\eta(j-\kappa_n)_+},\qquad
 d_j\max M\le C_\eta(1+j)(b_n+\kappa_n)
                         e^{\eta(j-\kappa_n)_+}.
\end{equation}
Fix one positive \(\eta\) to use the higher-moment assertion, and
let \(p>2\) be the resulting exponent. On the row \(\kappa_n\)
put
\[
 E_n=\{x:d_{\kappa_n}|f(x+\ii(H-d_{\kappa_n}))|\ge b_n/2\},
 \qquad p_n=\Prob(E_n).
\]
Holder's inequality gives
\[
 \frac34 b_n^2
 \le C(b_n+\kappa_n)^2p_n^{1-2/p}.
\]
Therefore
\begin{equation}\label{critical:eq:event-cost}
 p_n\ge c\left(\frac{b_n}{b_n+\kappa_n}\right)^{2p/(p-2)}
       =\kappa_n^{-o(1)}.
\end{equation}

We next select a normal level for the full moment. For any fixed
small \(\eps>0\), maximize
\(\mu(j)e^{-\eps|j-\kappa_n|}\) over \(j\ge k_0\).
Let \(k\) be a maximizer and \(B=\mu(k)\). The triangle
inequality gives
\begin{equation}\label{critical:eq:discount-normality}
 \mu(j)\le B e^{\eps|j-k|},\qquad B\ge b_n.
\end{equation}
Using \eqref{critical:eq:record-envelope} with rate \(\eps/2\)
also gives
\[
 |k-\kappa_n|\le\frac2\eps
         \log\left(\frac{C_\eps(b_n+\kappa_n)}{b_n}\right).
\]
For fixed \(\eps\), the right side is \(o(\log\kappa_n)\),
by \eqref{critical:eq:f-record}. A slow diagonal choice thus gives
\(\eps_n\downarrow0\) and maximizing levels \(k_n\) such that
\begin{equation}\label{critical:eq:small-shift}
 B_n^{\mathrm{new}}:=\mu(k_n)\ge b_n,
 \qquad |k_n-\kappa_n|=o(\log\kappa_n),
\end{equation}
while \eqref{critical:eq:discount-normality} holds with
\(\eps=\eps_n\). For the remainder of this proof write
\(B_n=B_n^{\mathrm{new}}\); the record value remains denoted by
\(b_n\).

Choose \(s_n\to\infty\), \(s_n=o(\log\kappa_n)\), so slowly
that it dominates both \(|k_n-\kappa_n|\) and \(\log(1/p_n)\).
Such a choice is possible by \eqref{critical:eq:event-cost} and
\eqref{critical:eq:small-shift}; for example take the geometric mean
of \(\log\kappa_n\) and
\(1+|k_n-\kappa_n|+\log(1/p_n)\).
Set \(k_{a,n}=k_n-s_n\), let \(\omega_n\) be uniform horizontal
measure conditioned on \(E_n\), and let
\(\nu_n=K_{d_{a,n}}*\omega_n\). Both \(\omega_n\) and
\(\nu_n\) have densities at most \(p_n^{-1}=e^{o(s_n)}\)
relative to uniform measure.

We check all ancestor bounds. Fix \(0<\beta<1/2\), a tile
enlargement, and \(0<\beta_0<\beta\). Primary normality and
tile averaging give the unweighted bound
\[
 \left\|\sup_{\mathcal T_A(j,x)}d_jM\right\|_2
 \le C_A B_n e^{\beta_0(k_n-j)}
 \quad(j\le k_{a,n}+3)
\]
for large \(n\). Passing to \(\nu_n\) costs at most
\(p_n^{-1/2}\). This is absorbed by
\(e^{(\beta-\beta_0)(k_n-j)}\), uniformly on that range because
\(k_n-j\ge s_n-3\). For the anchor, apply the \(P\) estimate
of \Cref{thm:critical-bootstrap} to the envelope
\[
 B_n\exp\bigl(\beta_0(k_n-j)_++a_0(j-k_n)_+\bigr)
\]
with any fixed \(a_0>0\). Primary normality supplies its hypotheses.
Conditioning costs \(p_n^{-1}\), which is absorbed by
\(e^{2(\beta-\beta_0)s_n}\). This proves
\eqref{critical:eq:prohibited-coarse} and
\eqref{critical:eq:prohibited-anchor}, with constants independent of
\(n\) for every fixed choice.

It remains to prove \eqref{critical:eq:prohibited-Q}. Fix a compact
set \(\mathcal K\) in the lower half-plane and a fixed \(A>1\).
For a center \(x_0\in E_n\), consider all target disks needed to
control \(q\) on \(x+\ii H+d_n\mathcal K\), where
\(|x-x_0|_{\mathrm{per}}\le A d_{a,n}\). One can connect the
initial point at level \(\kappa_n\) to these disks by first moving
to the ancestor row, moving horizontally there, and then moving to
the target depth. The number of depth-sized disks required is
\(O_A(1+s_n+|k_n-\kappa_n|)=o(\log\kappa_n)\).

On a fixed enlargement of this entire range,
\eqref{critical:eq:record-envelope} implies
\[
 \log|f|\le U_n^*:=\kappa_n+\log b_n+C\log\kappa_n.
\]
Here \(C\) can be fixed: all level changes are
\(o(\log\kappa_n)\), and
\(\log((b_n+\kappa_n)/b_n)=o(\log\kappa_n)\).
At the initial point in \(E_n\), the nonnegative superharmonic
deficit \(U_n^*-\log|f|\) is \(O(\log\kappa_n)\).
The disk-chain estimate \eqref{critical:eq:deficit-step} bounds its
average on every target disk by \(\kappa_n^{o(1)}\), uniformly in
the indicated centers. \Cref{lem:log-deficit}, with fixed inner
margins, therefore gives
\[
 \sup_{S\in\mathcal K}d_n
       |q(x+\ii H+d_nS)|^{1/2}\le\kappa_n^{o(1)}.
\]
Since \(B_n\ge b_n\ge\kappa_n^{1-o(1)}\), this is negligible
relative to \(B_n\).

Under \(\nu_n\), a center can be sampled by first choosing
\(x_0\) with law \(\omega_n\) and then adding the kernel
increment. The probability that its periodic distance exceeds
\(A d_{a,n}\) is at most \(C/A\), uniformly in \(n\).
For every fixed \(A\), the preceding estimate holds on the
complement. First let \(n\to\infty\), and then let \(A\to\infty\).
This proves \eqref{critical:eq:prohibited-Q}. We have constructed all
the data of \Cref{def:prohibited-sequence}, a contradiction.
\end{proof}

\begin{corollary}\label{cor:critical-polynomial-bounds}
Assume that no sequence in \Cref{def:prohibited-sequence} exists.
Then, for all \(k\ge k_0\),
\begin{equation}\label{critical:eq:polynomial-bounds}
 \mu(k)\le C(1+k),\qquad
 d_k\max_x M\le C(1+k)^2,\qquad
 d_k\max_x|P|\le C(1+k)^4.
\end{equation}
All fixed tile enlargements have the corresponding bounds.
\end{corollary}

\begin{proof}
By \Cref{critical:prop:power-saving} and
\eqref{critical:eq:q-first-moment},
\(\E d_jM\le C(1+j)\). At a target level \(K\), use the
nondecreasing envelope
\[
 N_K(j)=
 \begin{cases}
  1+j,&k_0\le j\le K,\\
  (1+K)e^{j-K},&j\ge K.
 \end{cases}
\]
Its envelope and first-moment constants are independent of \(K\),
and its exponential tail makes the bootstrap ratio finite. Apply
\Cref{thm:critical-bootstrap} at \(K\). Since \(N_K(K)=1+K\),
all three assertions follow. Fixed tile versions follow by applying
the same row estimates on their fixed ranges of depth ratios.
\end{proof}

\begin{lemma}\label{critical:lem:growth-charge}
Assume that no sequence in \Cref{def:prohibited-sequence} exists,
and fix \(0<\alpha<1/2\). Set
\[
 G_f(j)=\log\max_x|f(x+\ii(H-e^{-j}))|.
\]
There are constants \(C,c>0\) and a sufficiently large \(K_0\)
such that every \(\kappa\ge K_0\) with
\(\mu_f(\kappa)\ge\kappa^\alpha\) satisfies
\begin{equation}\label{critical:eq:growth-charge}
 G_f(\kappa+C\log\kappa)-G_f(\kappa-C\log\kappa)
 \ge c\mu_f(\kappa).
\end{equation}
\end{lemma}

\begin{proof}
The function \(G_f\) is nondecreasing: apply the maximum principle
between two opposite rows and use real symmetry. Suppose the claimed
constants do not exist. We may choose \(\kappa_n\to\infty\) and
\(C_n\to\infty\), with \(C_n\log\kappa_n=o(\kappa_n)\), such
that \(b_n:=\mu_f(\kappa_n)\ge\kappa_n^\alpha\) and
\begin{equation}\label{critical:eq:failed-charge}
 G_f(\kappa_n+C_n\log\kappa_n)
       -G_f(\kappa_n-C_n\log\kappa_n)=o(b_n).
\end{equation}
For clarity one can first set \(C_n=n\), require the increment to
be less than \(b_n/n\), and then choose the violating level larger
than any prescribed threshold.

Maximize \(\mu(j)e^{-\eps_n|j-\kappa_n|}\), where
\(\eps_n\downarrow0\) sufficiently slowly, and denote the
maximizer by \(k_n\), with \(B_n=\mu(k_n)\ge b_n\).
The maximizing argument in \eqref{critical:eq:discount-normality}
gives primary normality. The polynomial bound \(\mu(j)\le C(1+j)\)
shows that
\[
 |k_n-\kappa_n|=O(\eps_n^{-1}\log\kappa_n).
\]
Indeed, on the future side use
\((1+j)/(1+\kappa_n)\le e^{(j-\kappa_n)/(1+\kappa_n)}\)
and choose \(\eps_n\ge2/(1+\kappa_n)\); the past side is
immediate. We can therefore arrange
\[
 |k_n-\kappa_n|=o(C_n\log\kappa_n),\qquad
 \log\kappa_n+|k_n-\kappa_n|\ll s_n
       \ll C_n\log\kappa_n.
\]
Here both ratios denoted by \(\ll\) tend to zero. For example,
after a further slow choice of \(C_n\), one may use
\(\eps_n=C_n^{-1/3}\) and \(s_n=C_n^{2/3}\log\kappa_n\).
Set \(k_{a,n}=k_n-s_n\); then \(k_{a,n}\to\infty\).
Choose a horizontal maximizer \(x_{a,n}\) of \(|f|\) on the
ancestor row, take \(\omega_n\) to be its point mass, and put
\(\nu_n=K_{d_{a,n}}*\omega_n\).

For every fixed \(\beta>0\), the coarse and anchor hypotheses
follow directly from \eqref{critical:eq:polynomial-bounds}.
Indeed all required coarser tile suprema are \(O_A((1+k_n)^2)\),
whereas \(B_n e^{\beta(k_n-j)}\) is at least
\(\kappa_n^\alpha e^{\beta(s_n-3)}\) for
\(j\le k_{a,n}+3\). The latter dominates every fixed polynomial
in \(\kappa_n\), since \(s_n/\log\kappa_n\to\infty\).
Similarly the normalized ancestor bound for \(P\) is
\(O((1+k_{a,n})^4)\), and is dominated by
\(B_n^2e^{2\beta s_n}\). These are supremum estimates, so no
bound on the density of the point mass is needed.

Fix a compact target set \(\mathcal K\subset\{\Im S<0\}\),
and choose a fixed \(C'>0\) so that the row \(k_n+C'\) is above
all the target disks and their fixed enlargements. On the periodic
strip from the real row to that row, put
\[
 U_n(s)=G_f(k_n+C')-\log|f(s)|\ge0.
\]
At the chosen ancestor point, monotonicity and
\eqref{critical:eq:failed-charge} give
\begin{equation}\label{critical:eq:small-ancestor-deficit}
 U_n(x_{a,n}+\ii(H-d_{a,n}))
   =G_f(k_n+C')-G_f(k_{a,n})=o(b_n)=o(B_n).
\end{equation}
Both levels in this difference lie inside the interval in
\eqref{critical:eq:failed-charge}, for sufficiently large \(n\).

For any row in the fixed target range, apply positive
superharmonic comparison to \(U_n\) from the ancestor, using the
real bottom. The upper-boundary harmonic-measure law is comparable
to \(\nu_n\), because its depth is \(o(d_{a,n})\). The bottom
contribution is nonnegative. It follows that
\[
 \int U_n(x+\ii(H-d_j))\,\dd\nu_n(x)=o(B_n)
\]
uniformly over that fixed range of \(j-k_n\). The comparison can
be justified first for the superharmonic truncations
\(\min\{U_n,R\}\) and then by monotone convergence. Averaging
over the target disks has the same bound: their horizontal shifts
are \(O(d_n)\), and the density of \(\nu_n\) changes by bounded
factors under those shifts.

Apply \Cref{lem:log-deficit} on a finite covering of the target
compact set by disks with fixed inner margins. It yields
\[
 \int \sup_{S\in\mathcal K}
          d_n|q(x+\ii H+d_nS)|^{1/2}\,\dd\nu_n(x)
 \le C+o(B_n).
\]
Since \(B_n\to\infty\), Markov's inequality proves
\eqref{critical:eq:prohibited-Q}. Together with the preceding
normality and ancestor estimates this constructs a prohibited
sequence, the required contradiction.
\end{proof}

\begin{theorem}[Conditional angle-budget reduction]\label{thm:angle-budget-reduction}
If no sequence in \Cref{def:prohibited-sequence} exists, then for
some \(\eps>0\) and \(C<\infty\),
\begin{equation}\label{critical:eq:angle-budget}
 \int_{k_0}^K\mu_f(k)^2\,\dd k\le C(1+K)^{2-\eps}
 \qquad(K\ge k_0).
\end{equation}
In particular \Cref{thm:sequence-exclusion} will establish this
conclusion and \Cref{cor:critical-polynomial-bounds} in the critical
case.
\end{theorem}

\begin{proof}
Fix \(0<\alpha<1/2\) and split the sufficiently large levels into
the high set \(\mathcal A=\{k:\mu_f(k)\ge k^\alpha\}\) and its
complement. By \Cref{critical:lem:growth-charge},
\[
 \int_{\mathcal A\cap[K_0,K]}\mu_f(k)\,\dd k
 \le C\int_{K_0}^K
       [G_f(k+C\log k)-G_f(k-C\log k)]\,\dd k.
\]
The polynomial bound gives \(G_f(v)\le v+C\log(2+v)\).
Since \(G_f\) is nondecreasing, its positive Stieltjes measure can
be used to evaluate the last integral. Each point of that measure
is counted for a set of \(k\)'s of length at most \(2C\log K\).
Its total mass over the relevant interval, whose upper endpoint is
\(K+C\log K\), is \(O(K)\). The lower endpoint can be kept
fixed by increasing \(K_0\). Hence
\[
 \int_{\mathcal A\cap[K_0,K]}\mu_f(k)\,\dd k=O(K\log K).
\]
Multiplying by the bound from
\Cref{critical:prop:power-saving} gives
\[
 \int_{\mathcal A\cap[K_0,K]}\mu_f(k)^2\,\dd k
 \le C K^{2-\eps_1}\log K.
\]
On the complementary set the integral is at most
\(\int_{K_0}^K k^{2\alpha}\,\dd k=O(K^{1+2\alpha})\).
Both have a positive power saving from \(K^2\); decreasing that
power slightly absorbs the logarithm and the initial compact range.
This proves \eqref{critical:eq:angle-budget}.
\end{proof}
